\chapter{Conclusion}
\label{chap:conclusion}

This thesis asked whether post-COVID ECB monetary easing transmitted more strongly through housing-credit and financial-intermediation channels than through compensation-linked real-economy variables. The answer is yes, within the governed monthly reduced-form local-projection framework. The strongest evidence is not a broad asset-market result. It is a channel-specific hierarchy centered on house-purchase lending growth and supported by lending-condition variables.

\section{Housing-Credit Transmission}

The clearest and most persistent response channel is housing credit. House-purchase lending growth responds immediately, peaks at h6, and remains positive in accumulated terms through h24. This makes it the centerpiece variable and the core empirical contribution of the thesis.

Pure-new house-purchase loans provide a more volatile flow-based check. Their weaker evidence does not overturn the result, because stock-based credit growth and flow-based originations measure different objects. The former captures sustained financing conditions; the latter captures immediate origination behavior. The final housing-credit conclusion is therefore strong but precise: sustained house-purchase credit growth is the most visible response surface in the model.

\section{Financial Markets and Information Effects}

Financial-market responses are mixed and require careful interpretation. Lending spreads provide supporting intermediation evidence: mortgage spreads and NFC spreads are positive and cumulatively durable, with the mortgage spread especially stable. Broad credit quantities and ECB assets are weaker.

The real DAX response is negative at medium and long horizons. This result rules out a simple narrative that easing surprises uniformly raise asset prices. The better interpretation is information-effect sensitivity: ECB announcements can convey adverse macroeconomic information as well as accommodative policy news, and equity markets react to both. The DAX result is therefore a boundary on broad asset-price claims, not a reason to discard the housing-credit evidence.

\section{Compensation and Real-Economy Evidence}

Compensation-linked and real-economy responses are weaker and less consistent relative to the housing-credit block. The ECB Wage Tracker excluding one-offs moves at short and intermediate horizons but loses h24 cumulative exclusion. Employment expectations reverse at long horizons. The German industry wage-bill proxy prevents an absence claim, but it remains a narrower robustness measure.

The conclusion is comparative. The thesis finds stronger relative transmission through housing credit and selected intermediation variables than through compensation-linked proxies. It does not convert those comparisons into household incidence claims.

\section{Interpretation Boundary}

The results should be read as reduced-form evidence of heterogeneous monetary transmission. They support channel-specific persistence, relative transmission strength, housing-credit amplification, and timing-consistent financial intermediation. They do not support household welfare conclusions, inequality decomposition, redistribution claims, or a definitive verdict on household income incidence.

This boundary is not a weakness of the thesis. It is what makes the contribution credible. The model uses high-frequency ECB easing surprises and dynamic impulse responses to show where policy news remains most visible. The strongest answer is disciplined: post-COVID ECB easing shocks transmitted most clearly and persistently through housing-credit and selected financial-intermediation channels, while evidence for broad compensation-linked real-economy transmission was weaker and less consistent.

\section{Future Research}

Future research can extend the thesis in four directions. Bank-level and loan-level data would test whether the timing patterns documented here reflect lender heterogeneity, funding-cost pass-through, margin adjustment, or credit-supply behavior. Household-level credit, income, and balance-sheet data would allow incidence analysis across borrowers, renters, homeowners, income groups, and refinancing status. Regional housing data would show whether the housing-credit response is stronger where local housing constraints are already tight. A longer post-COVID sample would improve regime analysis and separate easing, tightening, balance-sheet normalization, and inflation-regime effects more convincingly.

The final position is therefore ambitious but bounded. This is a rigorous macro-financial transmission thesis examining whether ECB easing propagated more strongly through housing-credit and selected financial channels than through compensation-linked real-economy variables. The evidence supports that framing, and it does so without asking the model to identify more than the design permits.
